\documentclass[a0paper,portrait]{baposter}
\usepackage[utf8]{inputenc}
\usepackage{cmap}
\usepackage[T1]{fontenc}
\usepackage[english]{babel}
\usepackage{biblatex}
\addbibresource{references.bib}

\usepackage{url}
\usepackage{breakurl}

\usepackage{amsmath}
\usepackage{amsfonts}
\usepackage{amssymb}
\usepackage{amsthm}

\usepackage{graphicx}
\usepackage{tikz}
\usepackage{pgfplots}
\pgfplotsset{compat=newest,compat/show suggested version=false}

\usepackage{textcomp}
\usepackage{listings}

\usepackage{lmodern}
\usepackage{mathrsfs}

\usepackage{multirow}

\usepackage{unicamplogo}
\usepackage{imecclogo}

\definecolor{bordercol}{RGB}{40,40,40}
\definecolor{headercol1}{RGB}{186,215,230}
\definecolor{headercol2}{RGB}{80,80,80}
\definecolor{headerfontcol}{RGB}{0,0,0}
\definecolor{boxcolor}{RGB}{186,215,230}

\begin{document}
\begin{poster}
  {
    grid=false,
    columns=2,
    headerfont=\Large\sf\bf,
    background=shadeTB,
    borderColor=bordercol,
    headerColorOne=headercol1,
    headerColorTwo=headercol2,
    headerFontColor=headerfontcol,
    boxColorOne=boxcolor,
    bgColorOne=white,
    bgColorTwo=white,
    textborder=rectangle,
    headerborder=open,
    headershape=rounded,
    boxshade=plain
  }
  {
    \unicamplogol
  }
  {
    perprof-py \\ A Python Package for Performance Profile
  }
  {
    \begin{tabular}{ccc}
      Raniere Gaia Costa da Silva &
      Abel Soares Siqueira &
      Luiz-Rafael Santos \\
      \url{raniere@ime.unicamp.br} &
      \url{abel.s.siqueira@gmail.com} &
      \url{lrsantos11@gmail.com}
    \end{tabular}
  }
  {
    \imecclogo
  }

  \begin{posterbox}[column=0]{Introduction}
    Benchmarking optimization packages are very important in optimization
    field, not only because it is one of the way to compare solvers, but also
    to uncover deficiencies that could be overlooked. During benchmarking, one
    can obtain several informations, like CPU time, number of functions
    evaluations, number of iterations and so on. These informations, if
    presented as tables, can be difficult to be analyzed, due, for instance,
    to large amount of data. Therefore, researchers started testing tools to
    better process and understand this data.  One of the most widespread ways
    to do so is using Performance Profile graphics proposed by
    \textcite{Dolan2001}.
  \end{posterbox}

  \begin{posterbox}[column=0,below=auto]{Performance Profile}
    The performance profile for a solver is the cumulative distribution
    function for a performance metric, e.g, CPU time, number of functions
    evaluations, number of iterations and so on.

    Given a set $P$ of problems and a set $S$ of solvers. For each problem $p
    \in P$ and solver $s \in S$, we define $t_{ps}$ as the computing time
    required to solve problem $p$ by solver $s$ and
    \begin{align*}
      r_{ps} = \frac{t_{ps}}{\min\{t_{ps}: s \in S\}}
    \end{align*}
    as the performance ratio of solver $s$ for the problem $p$ when compared
    with the best performance by any solver on this problem.

    \textcite{Dolan2001} define
    \begin{align*}
      \rho_s(\tau) = \frac{| \{p \in P: r_{ps} \leq \tau\} |}{| P |}
    \end{align*}
    where $\rho_s(\tau)$ is the probability for solver $s \in S$ to solve one
    problem within a factor $\tau \in \mathbb{R}$ of the best performance
    ratio. For a given $\tau$, the best solver is the one with the higher
    value for $\rho_s(\tau)$.

    One of the advantages of using performance profiles is the easily graphic
    analyse, it be relatively insensitive to changes in results on a small
    number of problems and also largely unaffected by small changes in
    results over many problems.
  \end{posterbox}

  \begin{posterbox}[column=0,below=auto]{perprof-py}
    We implemented a free/open software that makes performance profile using data
    provided by user in a friendly manner. This software produces graphics in PDF
    using LaTeX with PGF/TikZ and pgfplots packages,
    in PNG using matplotlib and can also be easily extended to
    use with other plot library. The software is implemented in Python3 with
    support for internationalization and is available on
    \url{https://github.com/lpoo/perprof-py}.
  \end{posterbox}

  \begin{posterbox}[column=0,below=auto,height=bottom]{Output}
    \includegraphics[width=0.95\linewidth]{abc-semilog.pdf}
  \end{posterbox}

  \begin{posterbox}[column=1]{Features}
    Some of perprof-py features are:
    \begin{itemize}
      \item high quality output (in more than one format),
      \item lots of native options/filters,
      \item options to treat some naive problems (e.g. spend time equals to zero),
      \item native internationalization support.
    \end{itemize}

    \begin{center}
      \begin{tabular}{cc}
        \includegraphics[width=0.45\linewidth]{abc-bw-semilog.pdf} &
        \includegraphics[width=0.45\linewidth]{abc.pdf} \\
        \includegraphics[width=0.45\linewidth]{abc-100.pdf} &
        \includegraphics[width=0.45\linewidth]{abc-hs}
      \end{tabular}
    \end{center}

    Beside the previous features, the package can be easily used in a range
    of external tools that support Python.
  \end{posterbox}

  \begin{posterbox}[column=1,below=auto]{Input file}
    Each solver to be used in the benchmark must have a file like:

    \begin{lstlisting}
#Name Solver
Problem01 exit01 time01 objective01 primal01 dual01
Problem02 exit02 time02 objective02 primal02 dual02
    \end{lstlisting}

    Each file has 6 columns that means, in order:
    \begin{itemize}
      \item The name of the problem (\textbf{mandatory});
      \item Exit flag (\textbf{mandatory});
      \item Time spent (\textbf{mandatory});
      \item Objective function value (\textbf{optional}, to be used in future
        features);
      \item Primal infeasibility (\textbf{optional}, to be used in future
        features);
      \item Dual infeasibility (\textbf{optional}, to be used in future
        features).
    \end{itemize}
  \end{posterbox}

  \begin{posterbox}[column=1,below=auto]{Future works}
    We have plans to extend perprof-py adding more features, improve metadata
    support and also support data profiles as define by \textcite{More2009}.
  \end{posterbox}

  \begin{posterbox}[column=1,below=auto]{References}
    \printbibliography[heading=none]
  \end{posterbox}

  \begin{posterbox}[column=1,below=auto,height=bottom]{License}
    \begin{tabular}{cl}
      \multirow{2}{*}{\includegraphics[height=24pt]{figures/cc-by}} &
      This work is licensed under a \\
      & Creative Commons Attribution 3.0 Unported License.
    \end{tabular}
  \end{posterbox}
\end{poster}
\end{document}
